\chapter*{پیش‌گفتار}
\addcontentsline{toc}{chapter}{پیش‌گفتار}

\begin{multicols}{2}

این کتاب از یک پرسش ساده آغاز شد: چرا ما همدیگر را نمی‌فهمیم؟

در گفتگوهای روزمره، در شبکه‌های اجتماعی، در اخبار، همیشه از «چپ» و «راست» می‌شنویم. این واژه‌ها گاه به عنوان توهین به کار می‌روند، گاه به عنوان هویت، گاه به عنوان تحلیل. اما کمتر کسی می‌داند این مفاهیم از کجا آمده‌اند، چه تحولاتی داشته‌اند، و چرا هنوز مهم‌اند.

من سال‌ها است که با علاقه‌ای جدی تاریخ اندیشهٔ سیاسی را مطالعه می‌کنم. هرچه بیشتر می‌خوانم، بیشتر متوجه می‌شوم که اکثر مردم—حتی تحصیل‌کرده‌ها—درک روشنی از این مفاهیم ندارند. مارکسیسم، لیبرالیسم، محافظه‌کاری، فاشیسم، لیبرتارینیسم... همه این واژه‌ها به کار می‌روند، اما معنای دقیق‌شان چیست؟ چه تفاوتی با هم دارند؟ از کجا آمده‌اند؟

این کتاب تلاشی است برای پاسخ به این پرسش‌ها. من سعی کرده‌ام تاریخ طیف سیاسی را از انقلاب فرانسه تا امروز، به زبانی ساده اما دقیق روایت کنم. هدفم این نبوده که خواننده را به سمت هیچ جریانی ببرم؛ هدفم این بوده که ابزار فکری لازم برای درک و قضاوت مستقل را فراهم کنم.

\subsection*{این کتاب برای چه کسانی است؟}

این کتاب برای هر کسی نوشته شده که می‌خواهد جهان سیاست را بهتر بفهمد:

\begin{itemize}
    \item دانشجویانی که می‌خواهند پایه‌ای محکم در علوم سیاسی داشته باشند
    \item شهروندانی که می‌خواهند رأی آگاهانه‌تری بدهند
    \item روزنامه‌نگاران و تحلیلگرانی که با این مفاهیم کار می‌کنند
    \item هر کسی که از بحث‌های سیاسی سردرگم شده و می‌خواهد نقشه‌ای داشته باشد
\end{itemize}

پیش‌نیاز خاصی لازم نیست. من فرض کرده‌ام که خواننده تحصیل‌کرده است اما لزوماً متخصص نیست.

\subsection*{چرا بی‌طرفی؟}

یکی از چالش‌های نوشتن این کتاب، حفظ بی‌طرفی بود. هر نویسنده‌ای دیدگاهی دارد، و من هم استثنا نیستم. اما تلاش کرده‌ام این دیدگاه را بر کتاب تحمیل نکنم.

چرا؟ چون هدف این کتاب آموزش است، نه تبلیغ. می‌خواهم خواننده پس از مطالعه بتواند خودش قضاوت کند. نمی‌خواهم جای خواننده فکر کنم.

البته بی‌طرفی کامل ممکن نیست. انتخاب موضوعات، لحن نوشتار، ترتیب فصل‌ها—همه اینها انتخاب‌هایی هستند که نویسنده می‌کند. اما تلاش کرده‌ام هر جریان را با انصاف معرفی کنم: هم استدلال‌هایش را بگویم، هم نقدهایی که به آن وارد شده.

\subsection*{ساختار کتاب}

کتاب در سه بخش اصلی سازماندهی شده:

\textbf{بخش اول} (فصل‌های ۱-۵) مبانی را می‌گذارد: طیف سیاسی چیست؟ از کجا آمده؟ و چپ کلاسیک (سوسیالیسم، مارکسیسم) چه می‌گوید؟

\textbf{بخش دوم} (فصل‌های ۶-۱۳) به بررسی جریان‌های مختلف می‌پردازد: سوسیال دموکراسی، آنارشیسم، محافظه‌کاری، نئولیبرالیسم، لیبرتارینیسم، ناسیونالیسم، و راست نو.

\textbf{بخش سوم} (فصل‌های ۱۴-۱۷) به میانه‌روی، تجربه‌های ترکیبی، جهان غیرغربی، و آینده می‌پردازد.

\textbf{ضمائم} شامل تایم‌لاین، شخصیت‌نامه، واژه‌نامه، و کتاب‌شناسی است.

\subsection*{سپاسگزاری}

از همهٔ کسانی که در نوشتن این کتاب یاری‌ام رساندند سپاسگزارم: از دوستانی که نسخه‌های اولیه را خواندند و نظر دادند، از خانواده‌ام که صبورانه تحمل کردند، و از خوانندگانی که با پرسش‌ها و نقدهایشان مرا به تفکر بیشتر واداشتند.

\vspace{5mm}

\begin{flushright}
مهدی سالم\\
ریچموندهیل\\
زمستان ۲۰۲۴
\end{flushright}

\end{multicols}